\section{Introduction}
\label{sec:intro}

Move is a smart-contract language originally developed for the Libra/Diem blockchain at Meta~\cite{libra_blockchain_white,OLD_MOVE_LANG,blackshear2020resources} and now deployed at scale on Aptos~\cite{MoveOnAptosBook}, a high-throughput Layer-1 platform that settles transactions for hundreds of millions of dollars in user assets on chain.
A defining trait of the language is its resource-oriented type discipline: digital assets are represented as Move \emph{resources} whose ownership and movement are tracked by linear-style abilities (|copy|, |drop|, |store|, |key|), ruling out accidental duplication or destruction at compile time.
The Aptos platform builds on this foundation with a system of objects, parallel execution, and a rich on-chain framework written in Move itself~\cite{AptosFrameworkBook,ParkZGXGCLC24}.

Beyond its type discipline, Move has been designed from the outset to admit \emph{offline formal verification}: the Move Prover translates Move programs and their specifications, written in the Move Specification Language, into verification conditions discharged by SMT solvers~\cite{DillGPQXZ22,GrieskampEtAl26HO}.
These techniques are powerful but inherently offline: they apply prior to deployment, depend on hand-written specifications, and produce no signal once the code is running on chain.

Closer to the execution path, the Move VM also performs \emph{static bytecode verification} every time a module is loaded, enforcing well-formedness, typing, and Move's reference discipline --- a Rust-like~\cite{rust,rustbelt} aliasing and lifetime regime detailed in Section~\ref{sec:move}.
The reference analysis in particular is intricate: the verifier maintains a borrow graph that conservatively approximates the aliasing relationships entailed by control- and data-flow.
Implementing this correctly is a delicate engineering task, and the Aptos bug bounty program has surfaced verifier bugs in the past.
In a security-oriented architecture, the cost of trusting a single intricate static analysis as the sole gate is too high.

The Aptos runtime therefore carries a second layer of \emph{runtime safety checks} that re-establish the critical invariants during execution.
The checks described in this paper fall into three groups.
\emph{Type checks} verify that the dynamic type of every value matches the static type the bytecode expects to consume.
\emph{Ability checks} verify that operations such as copy, drop, or storage in global state are only performed on values whose declared abilities permit them.
\emph{Reference checks} verify, at runtime, the same dangling- and aliasing-freedom properties that the static verifier is supposed to enforce.

Since these checks are costly, the runtime supports \emph{trusted code} that waives them under certain conditions, and an \emph{asynchronous} checking mode that leverages idle workers in Block-STM~\cite{BlockSTM}, Aptos's speculative parallel executor.

\tengzhang{I think we should explicitly state the contribution: 1) propose a second layer of in-depth-defense for blockchain which potentially holds millions of dollars, 2) present a set of dynamic checks that are necessary for the second layer to be effective, 3) present a mechanism for async check by leveraging block-STM, which is originally designed for throughput.}
